\section{Raw oracle bounds for globally Lipschitz activations}
\label{sec:general-lipschitz-raw}

Fix a scalar activation $\phi$ and a constant $A\geq1$ such that
$\operatorname{Lip}(\phi)\leq A$ and $|\phi(0)|\leq A$.
The activation need not be bounded, homogeneous, odd, or zero at the
origin.  In particular the argument covers tanh, ReLU, GELU, sigmoid,
and softplus, with their appropriate fixed Lipschitz constants.

There are $H\geq1$ hidden layers of width $N$, a scalar linear readout,
and a first-layer constant input.  Internal layers have no biases.
Inputs belong to $[-1,1]^d$, labels to $[-1,1]$, and their joint IID
law is arbitrary.  Write $\tau\in(0,1]$ for the actual Gaussian scale:
first-layer, internal, and readout coordinate variances are
$\tau^2/d$, $\tau^2/N$, and $\tau^2/N^2$.
Put $\omega=N/\tau^2$.  The likelihood variance is $n\lambda/\omega$,
so the actual posterior is
\[
 d\Pi=Z^{-1}\exp[-\omega\widehat R_S(z_w)/(2\lambda)]\,d\mu.
\]
Taking $\tau^2=N^{1-\gamma}$ gives the original prefactor-one priors
at every $\gamma\geq1$.  No scale mixture appears in this learner.

\subsection{A teacher cylinder without bounded activations}

Let $g$ be a width-$m$ teacher, with all coefficients at most $B\geq1$
in absolute value.  Define the weighted Gaussian energy
\[
 Q_N(w)=\frac dN\|W^1\|_F^2+
             \sum_{\ell=2}^H\|W^\ell\|_F^2+N\|v\|_2^2.
\]
The Gaussian density has exponent $-\omega Q_N(w)/2$.
To represent the teacher at any $N\geq m$, partition each hidden
layer into clone groups of sizes
$k_j\in\{\lfloor N/m\rfloor,\lceil N/m\rceil\}$.
Copy each teacher first-layer row to its group. For an internal
teacher edge from neuron $i$ to neuron $j$, assign its weight divided
by $k_i$ to every edge between the corresponding clone groups.
Divide each teacher readout weight by its group size in the same way.
The sums entering every clone equal the teacher sums, so this network
computes exactly $g$, for any activation.

The first-layer contribution is at most $d(d+1)B^2$.
Since $k_j/k_i\leq2$, each internal-layer contribution is at most
$2m^2B^2$. Since $N/k_i\leq2m$, the readout contribution is at most
$2m^2B^2$. Thus the clone $w_*$ satisfies
\[
 Q_N(w_*)\leq E:=d(d+1)B^2+2Hm^2B^2.
\]
Every learner neuron belongs to a nonempty clone group, so this
construction does not require $\phi(0)=0$.
Define
\begin{gather*}
 b_0=2\sqrt{2[2\log9+\log(4(H+1))]},\qquad
 T_g=A[3+2mB+B(d+1)+b_0],\\
 P_g=1+(\sqrt2mB+1)A(d+1+H\sqrt m)T_g^{H-1}+T_g^H,\\
 D_g=d+2+(H-1)m,\qquad
 \eta_\rho=\min\{1/2,\rho/(2P_g)\},\qquad c_0=3/8,\\
 \mathcal A_g(\rho)=E/2+D_g\log(1/(c_0\eta_\rho))+\log2.
\end{gather*}
Let $U\geq\|g\|_\infty$ be any specified output bound;
one may always take $U=mBT_g^H$, and for tanh one may take $U=mB$.

\begin{lemma}[A finite-width teacher neighborhood]
\label{lem:general-activation-cylinder}
For every $N\geq m$, $0<\tau\leq1$, and $0<\rho\leq1$,
\[
 -\log\mu\{\|z_w-g\|_\infty<\rho\}
 \leq\omega E/2+ND_g\log(1/(c_0\eta_\rho))+\log2
 \leq\omega\mathcal A_g(\rho).
\]
\end{lemma}
\begin{proof}
Let $w_*$ be the balanced clone.  At each hidden layer let $P$ project
onto the $m$-dimensional subspace of vectors constant on clone groups.
In standardized independent Gaussian coordinates, require absolute
value at most $\eta_\rho$ for all first-layer and readout coordinates,
and for every internal row's projections onto an orthonormal basis of
this group subspace.  There are $ND_g$ independent scalar constraints,
whose joint probability is at least $(c_0\eta_\rho)^{ND_g}$.

Conditioning on these intervals preserves independent centered
unit-subgaussian coordinates in the corresponding orthonormal bases.
Two sphere nets, each of size at most $9^N$, give for an internal
standardized matrix $Z$
\[
 \Pr\{\|Z\|_{\rm op}>b_0\sqrt N\mid\text{cylinder}\}
 \leq2\exp[(2\log9-b_0^2/8)N].
\]
A union bound over internal layers leaves a centrally symmetric event
of probability at least one half the cylinder probability.

Translate this event by $w_*$.  The internal perturbation has operator
norm at most $b_0$, while its restriction to the teacher's group
subspace has operator norm at most $\eta_\rho\sqrt m$.
The clone internal operator norm is at most $\sqrt2mB$.
Teacher hidden coordinates are at most $T_g^\ell$ at layer $\ell$,
because $|\phi(s)|\leq A(1+|s|)$.
For normalized hidden errors $e_\ell=\|h^\ell-h_*^\ell\|_2/\sqrt N$,
Lipschitz propagation gives
\[
 e_1\leq A\eta_\rho(d+1),\qquad
 e_\ell\leq A(\sqrt2mB+b_0)e_{\ell-1}
                +A\eta_\rho\sqrt m\,T_g^{\ell-1}.
\]
Consequently
$e_H\leq\eta_\rho A(d+1+H\sqrt m)T_g^{H-1}$.
The clone readout has $\sqrt N\|v_*\|_2\leq\sqrt2mB$, and its
perturbation has $\sqrt N\|v-v_*\|_2\leq\eta_\rho$.
Thus the output error is at most
$(\sqrt2mB+1)e_H+\eta_\rho T_g^H<\rho$.

For a centered Gaussian measure and a centrally symmetric event,
translation by a vector of squared standardized length $r$ decreases
its probability by at most the factor $e^{-r/2}$; this follows by
pairing opposite points and applying Jensen to the Gaussian density
ratio.  Here $r\leq\omega E$.  This proves the first bound; $N\leq\omega$
and $1\leq\omega$ prove the second.
\end{proof}

\subsection{Width-uniform prior moments and norm tails}

Define
\[
 K(w)=\max\{128,\|W^1\|_F^2/N,
          \|W^2\|_{\rm op}^2,\ldots,\|W^H\|_{\rm op}^2,
          N\|v\|_2^2\}.
\]
Gaussian chi-square tails for the first layer and readout, and two
sphere nets for internal layers, imply for $u\geq128$
\begin{equation}
 \mu\{K>u\}\leq(H+1)e^{-\omega u/16}.
 \label{eq:general-activation-norm-tail}
\end{equation}
For example the internal estimate is
$2\exp[2N\log9-\omega u/8]\leq e^{-\omega u/16}$.
The first-layer term in $K$ is $(\tau^2/(dN))\chi^2_{N(d+1)}$;
the readout term is $(\tau^2/N)\chi^2_N$, so the same weaker bound
applies to them.
Tail integration yields
\begin{equation}
 \log\int e^{\omega K/32}d\mu\leq\omega B_H,
 \qquad B_H=4+\log(H+2).
 \label{eq:general-activation-norm-mgf}
\end{equation}

There are also width-uniform raw-output moments.  Let
$S_\ell=\|h^\ell(x)\|_2^2/N$, and use $S_0=\|\bar x\|_2^2/d\leq2$
for the first conditional Gaussian field.  Conditional on previous
layers, each new field has variance $\tau^2S_{\ell-1}$.
The bounds $|\phi(s)|\leq A(1+|s|)$ and Gaussian second and fourth
moments give
\[
 \mathbb E S_\ell\leq2(4A^2)^\ell,
 \qquad \mathbb E S_\ell^2\leq4(32A^4)^\ell.
\]
For the second assertion, Jensen applied to the average of squared
coordinates and $|\phi(s)|^4\leq8A^4(1+|s|^4)$ suffice; these bounds
also apply to the first layer with the stated $S_0$.
The centered Gaussian readout therefore implies, uniformly in $x$,
\[
 \mathbb E_\mu z_w(x)^2\leq 2(4A^2)^H/N,
 \qquad \mathbb E_\mu z_w(x)^4\leq12(32A^4)^H/N^2.
\]
Put
\[
 R_0=1+2(4A^2)^H,\qquad V_0=8+96(32A^4)^H.
\]
Tonelli and Jensen give, for every sample and every data law under
consideration,
\begin{equation}
 \int\widehat R_S(z_w)d\mu\leq R_0,\qquad
 \int R(z_w)d\mu\leq R_0,\qquad
 \int R(z_w)^2d\mu\leq V_0,
 \qquad Z\geq e^{-\omega R_0/(2\lambda)}.
 \label{eq:general-activation-prior-moments}
\end{equation}

\subsection{A common function cover and raw statistical transfer}

Write $L=H+1$.  On $K(w)\leq k$, Lipschitz propagation gives
\[
 \operatorname{Lip}(z_w)\leq A^H k^{L/2},\qquad
 \|z_w\|_\infty\leq C_{d,H,A}k^{L/2}.
\]
For the amplitude bound, normalized hidden norms satisfy
$a_\ell\leq A(1+\sqrt k\,a_{\ell-1})$ starting from
$a_0\leq\sqrt{d+1}$, and the readout multiplies by at most $\sqrt k$.
Thus this bound includes nonzero $\phi(0)$.
An input grid of mesh proportional to $a/k^{L/2}$, followed by rounding
values to multiples of $a$, provides a deterministic finite cover
$\mathcal C_{k,a}$ in uniform norm, common to all widths and scales,
with centers bounded by $C_{d,H,A}k^{L/2}+1$ and
\begin{equation}
 \log|\mathcal C_{k,a}|
 \leq C_{d,H,A} k^{dL/2}a^{-d}
                  \log\frac{C_{d,H,A}k}{a}
 \qquad(k\geq128,\ 0<a<1).
 \label{eq:general-activation-cover}
\end{equation}
Nearest-grid-point functions are sufficient; the cover centers need
not themselves be networks.

\begin{theorem}[Fixed-dimensional raw oracle inequality]
\label{thm:general-activation-fixed-d}
Fix $d,H,A$ as above and set
\[
 p_{d,H}=\frac{(H+1)(d+2)}2.
\]
There is a finite computable constant $C_{d,H,A}$ such that, whenever
$0<\lambda\leq1$, $0<\delta<1$, $n\geq2$, and
\begin{equation}
 n\geq C_{d,H,A}\lambda^{-p_{d,H}}
             \log^6\!\frac{2n}{\lambda\delta},
 \label{eq:general-activation-budget}
\end{equation}
for every fixed width-$m$, weight-$B$ teacher $g$, with probability
at least $1-\delta$ simultaneously at every $N\geq m$ and every
$0<\tau\leq1$,
\begin{align}
 \int R(z_w)d\Pi
 \leq\;&R(g)+2\lambda\mathcal A_g(\rho)+C_{d,H,A}\lambda\notag\\
 &+(U+1)\sqrt{\frac{2R(g)\log(2/\delta)}n}
       +\frac{2(U+1)^2\log(2/\delta)}{3n},
 \qquad \rho=\frac\lambda{4(U+1)}.
 \label{eq:general-activation-oracle}
\end{align}
The same bound holds for the ordinary raw posterior mean.
In particular it holds on one event for every original exponent
$\gamma\geq1$ and every eligible learner width, without changing the
prior prefactor.  At fixed teacher resources the rate is
$\widetilde O(n^{-2/((H+1)(d+2))})$.
\end{theorem}

\begin{proof}
Take $k_j=128\cdot2^j$ up to the least $J$ with
$k_J\geq(16R_0+128)/\lambda$, and use cover tolerances
$a_j=c_{H,A}\lambda k_j$ with a sufficiently small fixed positive
$c_{H,A}$.  In particular all tolerances are below one.
Apply Lemma~\ref{lem:finite-cover-bernstein} to these finite classes,
dividing failure probability $\delta/2$ among them. Their loss ranges
are at most $C_{d,H,A}k_j^L$.
When $\widehat R_S(z_w)\leq R_0+8$, an approximating center has
empirical risk bounded by a constant depending only on $H,A$.
The finite-cover Bernstein bound therefore makes its population-minus-empirical
risk a sufficiently small multiple of $\lambda k_j$, provided
\[
 n\geq C_{d,H,A}\log^6\!\frac{2n}{\lambda\delta}
    \max_{j\leq J}
       \frac{k_j^{L+dL/2}}{(\lambda k_j)^{d+2}}.
\]
The displayed maximum is at most $C_{H,A}\lambda^{-p_{d,H}}$:
$p_{d,H}\geq d+2$ and $k_j\leq C_{H,A}/\lambda$.
This verifies the condition from~\eqref{eq:general-activation-budget}.
Minkowski transfers population and empirical square-root risks between
a network and its center.  Their risks remain bounded by constants on
this region, so taking $c_{H,A}$ small enough gives
\[
 R(z_w)\leq\widehat R_S(z_w)+\lambda K(w)/16
 \quad\text{when }K(w)\leq k_J,\qquad
                       \widehat R_S(z_w)\leq R_0+8.
\]
Here select a dyadic shell $K(w)\leq k_j\leq2K(w)$ and absorb that
factor into the constants.  This one sample event covers all widths
and all scales, since the finite classes did.

The two omitted regions have small weighted posterior risk.
Using~\eqref{eq:general-activation-prior-moments},
\[
 \int R(z_w)\mathbf1_{\{\widehat R_S>R_0+8\}}d\Pi
 \leq R_0e^{-4\omega/\lambda}.
\]
Cauchy--Schwarz,~\eqref{eq:general-activation-norm-tail}, and the
normalizer bound give
\[
 \int R(z_w)\mathbf1_{\{K>k_J\}}d\Pi
 \leq\sqrt{V_0(H+1)}
     e^{\omega R_0/(2\lambda)-\omega k_J/32}
 \leq\sqrt{V_0(H+1)}e^{-4\omega/\lambda}.
\]
These deterministic estimates are simultaneous in all scales and
widths.  From~\eqref{eq:general-activation-norm-mgf} and posterior
variational optimality, for every comparison law $Q$,
\[
 \int[\widehat R_S+\lambda K/16]d\Pi
 \leq\int\widehat R_SdQ+
          \frac{2\lambda}\omega\operatorname{KL}(Q\Vert\mu)
          +2\lambda B_H.
\]
Choose $Q$ to be the prior conditioned on the teacher cylinder of
radius $\rho$.  Its empirical risk is at most
$\widehat R_S(g)+2\rho(U+1)+\rho^2\leq\widehat R_S(g)+\lambda$,
and Lemma~\ref{lem:general-activation-cylinder} bounds its entropy cost
by $2\lambda\mathcal A_g(\rho)$.
Integrate the deviation bound, add the two tails, and use
$e^{-4/\lambda}\leq\lambda$ to obtain the claimed bound with
$\widehat R_S(g)$ in place of the final teacher-risk terms.
Bernstein's inequality for this fixed teacher supplies those terms
with remaining failure probability $\delta/2$.

Finally the prior fourth moment is finite and the posterior density
relative to the prior is bounded by $1/Z$.  The raw posterior mean
therefore exists, and Jensen proves its risk bound.
\end{proof}

\subsection{Competing with a growing class}

Fix $d,H,A,\delta$, and let $p=p_{d,H}$. Choose the likelihood parameter
of order $\widetilde O(n^{-1/p})$, with a sufficiently large logarithmic
factor to meet \eqref{eq:general-activation-budget}. The same learner
then competes with the population-risk minimizer among teachers whose
width and weight caps are at most $n^\alpha$, whenever
\[
 \alpha<\frac1{4p}=\frac1{2(H+1)(d+2)}.
\]
Its oracle excess is $\widetilde O(n^{-1/p+4\alpha})$.
Indeed the teacher energy is $O(n^{4\alpha})$, the remaining cylinder
cost is $O(n^\alpha\log n)$, and one may take
$U=O(n^{2\alpha(H+1)})$. The minimizer has risk at most one because
zero is an available prediction. Its Bernstein terms therefore vanish
faster than the displayed excess, using $p\geq H+1$ and the stated
range of $\alpha$.

With a fixed weight cap and width cap $n^\alpha$, the sufficient range
and excess are instead
\[
 \alpha<\frac1{2p},\qquad
 \widetilde O(n^{-1/p+2\alpha}).
\]
Compactness of each capped parameter class and continuity of the risk
give a minimizer chosen before sampling. The Gaussian likelihood
calibration depends on the declared resources and confidence level,
not on that minimizer or on the unknown data distribution.

\subsection{A faster bound when averaging clipped outputs}

The likelihood can remain exactly the same while the prediction
convention clips individual network outputs before averaging.  Write
$c(w,x)=\operatorname{clip}_{[-1,1]}z_w(x)$ and
$\widehat f_S^{\rm clip}(x)=\int c(w,x)d\Pi(w)$.
This is generally different from clipping the ordinary raw posterior
mean, so the following refinement specifically concerns the displayed
predictor and individual clipped posterior draws.

\begin{corollary}[Clipped-output oracle bound]
\label{cor:general-activation-clipped}
Theorem~\ref{thm:general-activation-fixed-d} also bounds both
\[
 \int R(c(w,\cdot))d\Pi
 \quad\text{and}\quad R(\widehat f_S^{\rm clip})
\]
under the weaker sample condition
\[
 n\geq C_{d,H,A}\lambda^{-p_{d,H}^{\rm clip}}
          \log^6\!\frac{2n}{\lambda\delta},\qquad
 p_{d,H}^{\rm clip}=\max\{d+2,d(H+1)/2\}.
\]
The actual posterior still uses the raw squared-loss likelihood and
the same single Gaussian prior.
\end{corollary}
\begin{proof}
Clipping preserves the Lipschitz constant and gives output amplitude
one.  The uniform cover construction therefore has the same power
$k^{d(H+1)/2}a^{-d}$ in its logarithmic cardinality, but now a constant
loss range.  Lemma~\ref{lem:finite-cover-bernstein} requires only
\[
 n\geq C_{d,H,A}\log^6\!\frac{2n}{\lambda\delta}
     \max_j\frac{k_j^{d(H+1)/2}}{(\lambda k_j)^{d+2}}.
\]
Since $128\leq k_j\leq C_{H,A}/\lambda$, this maximum is bounded
by a constant times $\lambda^{-\max\{d+2,d(H+1)/2\}}$.
The empirical clipped risk is at most four, so no empirical-risk
truncation is needed.  The resulting common sample event gives
\[
 R(c(w,\cdot))\leq\widehat R_S(c(w,\cdot))+\lambda K(w)/16
             \leq\widehat R_S(z_w)+\lambda K(w)/16
 \qquad(K(w)\leq k_J).
\]
The omitted norm tail is controlled by the same weighted raw-risk
estimate, since clipping cannot increase loss against $|Y|\leq1$.
The entropy and teacher comparisons are unchanged.  Integrating and
then applying Jensen proves the two stated conclusions.
\end{proof}
